\section{Introduction}

Manipulator design establishes fundamental performance boundaries through the interplay of mechanical structure, actuation system, and control architecture.
At the hardware level, these capabilities manifest as explicit physical constraints: joint angle limits that define the reachable workspace, maximum velocities bounded by motor specifications, and torque limits determined by the drivetrain capacity.
However, a crucial distinction exists between these hardware limitations and the operational capacities specified by manufacturers for general-purpose deployment.
While hardware limits are absolute physical constraints, nominal specifications---like payload capacity and effective workspace boundaries---are \textit{derived} conservatively by considering worst-case scenarios across possible configurations and trajectories.
In practice, system integrators and end-users are locked into treating these derived operational capacities as hard limits rather than contextual capabilities, constrained by warranty requirements to design processes well within these conservative bounds.
This well-established approach of \textit{over-provisioning} hardware to ensure ample safety buffers has system-wide cascading effects---from the need to procure oversized actuators and power supplies, to reinforcing mounting structures for added load tolerance, to expanding facility requirements to accommodate larger equipment.
For instance, when an application requires handling a $35$kg payload, end-users are forced to select a $50$kg-rated robot over a $30$kg-rated system, despite the latter potentially being capable of safely executing the task within broad regions of its operational space.

Importantly, it is possible to overcome the issue of hardware over-provisioning through software-based optimizations, which allow to expand the safe operational envelope of a fixed/given robot embodiment and enable it to perform beyond nominal payload specifications.
This requires unified handling of both geometric (i.e., kinematic) and dynamic constraints, as payload characteristics interact with joint configurations and higher-order derivatives to determine feasible motions.
Even if they have not solved this specific problem, prior approaches in this space have investigated the challenge from multiple perspectives: geometric planners with dynamics post-processing \cite{orthey2023sampling,berscheid2021jerk}, kinodynamic planning \cite{shome2021asymptotically}, optimization-based methods \cite{ichnowski2022gomp}, and control-based frameworks \cite{arrizabalaga2024slosh}. However, these methods face inherent trade-offs between design complexity, constraint handling, and computational efficiency; these become particularly critical when operating near system limits where the feasible solution space becomes highly constrained.
These limitations highlight the need for a method that balances constraint satisfaction with
generalizability---i.e., the ability to adapt to varying payloads, environmental conditions, and task requirements while maintaining consistent performance.

Of note, diffusion models \cite{chi2023diffusionpolicy,sridhar2024nomad} present a compelling opportunity
address the ``curse of dimensionality'' inherent in trajectory planning with simultaneous kinematic and dynamic constraint satisfaction.
\textit{By learning from successful trajectories generated by multiple (imperfect) planning methods}, these models can extract various homotopic solution classes that overcome the limitations of individual approaches.
This distills successful solutions from diverse planning approaches into a unified representation that generalizes across the entire feasible space, including regions where individual planners fail.
The inherent stochasticity in the denoising process enables the exploration of diverse, multimodal trajectories \cite{chi2023diffusionpolicy}, while their capacity to model high-dimensional distributions allows incorporation of rich contextual information such as sensor data and problem constraints.
However, existing diffusion-based approaches have been confined to position-controlled systems \cite{carvalho2023motion,saha2024edmp}, leaving the fundamental challenge of integrating both kinematic and dynamic constraints unaddressed.

In this work, we propose a method that directly tackles this problem by jointly learning valid distributions of positions, velocities, and accelerations that satisfy both constraint types simultaneously---moving beyond prior work that only denoises in position space.
Through investigation of payload encoding strategies, we demonstrate that diffusion models effectively learn the complex relationships between joint states and payload-dependent dynamics, \uline{enabling constant-time generation ($\approx 10$ms) of feasible trajectories that preserve $67.6\%$ workspace accessibility at $3\times$ nominal payload capacity}---without explicit constraint checking at runtime. This demonstrates the importance of dynamics-informed generative models in expanding the operational boundaries of robotic systems.